% ---
% id: foundations/def:binary64
% kind: def
% title: "Binary64 numbers, declared values and dyadic lattices"
% summary: >-
%   A declared value is a binary64 number and denotes exactly the dyadic
%   rational it is, with no exact object behind it. A binary64 addition is
%   exact precisely when the true sum is a binary64 number, and a lattice value
%   at k is an integer multiple of two to the power −k of modulus at most one.
% prominence: standard
% status: implemented
% uses: []
% assumes: []
% refs: []
% provenance: newly-derived
% implements: certiqs_key.foundations.exact:scalar_is_binary64
% lean:
% ---

\begin{definition}\label{def:binary64}
\begin{enumerate}[label=(\roman*),leftmargin=2.4em]
\item Let $\mathbb{B}$ be the set of real numbers $m\,2^{q}$ with
  $m,q\in\mathbb{Z}$, $|m|<2^{53}$ and $-1074\le q\le971$, and let
  $\mathbb{B}_{\mathbb{C}}\coloneqq\{x+iy:\ x,y\in\mathbb{B}\}$. The elements
  of $\mathbb{B}$ are the \term[the finite values of the IEEE 754 binary64
  format, each a dyadic rational]{binary64 numbers}, the finite values of the
  IEEE~754 binary64 format. Every binary64 number is a dyadic rational, an
  element of $\mathbb{Z}[\tfrac12]$.
\item For a real number $x$ with $|x|<2^{1024}-2^{970}$, $\mathrm{fl}(x)$ is
  $x$ rounded to the nearest element of $\mathbb{B}$, a tie going to the
  element whose $m$ is even in its representation of (i) with the least $q$.
  For $|x|\ge2^{1024}-2^{970}$, $\mathrm{fl}(x)\coloneqq+\infty$ if $x>0$ and
  $-\infty$ if $x<0$; the threshold is the midpoint between the largest
  element $(2^{53}-1)\,2^{971}$ of $\mathbb{B}$ and $2^{1024}$, which is where
  IEEE~754 rounding to nearest overflows. Further,
  $\mathrm{fl}(\pm\infty+y)\coloneqq\pm\infty$ for $y\in\mathbb{B}$. The
  \term[fl(a + b), the exact sum of a and b rounded to the nearest binary64
  number, ties to even, or to an infinity on overflow]{binary64 sum} of $a$
  and $b$ is $\mathrm{fl}(a+b)$. On
  $\mathbb{B}_{\mathbb{C}}$, and on matrices over it, $\mathrm{fl}$ acts on the
  real and on the imaginary part of every entry separately. Two properties of
  $\mathrm{fl}$ are all that is used of it: on real input its value lies in
  $\mathbb{B}\cup\{\pm\infty\}$, and $\mathrm{fl}(x)=x$ for $x\in\mathbb{B}$.
  The binary64 sum of $a$ and $b$ is \emph{exact} when it equals $a+b$; by
  those two properties it is exact if and only if the real and the imaginary
  part of every entry of $a+b$ lie in $\mathbb{B}$.
\item A \term[a binary64 number stated as a datum, which denotes exactly that
  dyadic rational and approximates nothing]{declared value} is a binary64
  number stated as a datum. It
  denotes that dyadic rational and nothing else: there is no exact object
  behind it of which it is an approximation, and every predicate stated of
  declared values is a predicate of those rationals. A declared complex value
  is an element of $\mathbb{B}_{\mathbb{C}}$, and a declared matrix is a
  matrix of declared complex values.
\item A \term[an integer k with 1 ≤ k ≤ 53]{lattice exponent} is an integer
  $k$ with $1\le k\le53$. A real number $x$ is a \term[at a lattice exponent
  k, a real number that is an integer multiple of two to the power −k and has
  modulus at most one; a complex number or a matrix is one when the real and
  the imaginary part of every entry are]{lattice value} at $k$ when
  $2^{k}x\in\mathbb{Z}$ and
  $|x|\le1$; a complex number, or a matrix, is one when the real and the
  imaginary part of every entry are.
\end{enumerate}
\end{definition}
